\section{Experimental design}
\label{sec:design}

\begin{figure}[H]
\centering
\includegraphics[width=0.86\linewidth,height=0.42\textheight,keepaspectratio]{experimental_design.png}
\caption{Design schematic. The diagram shows the axis set as originally specified, before
the 2026-06-10 trim removed two prompt levels and the \code{plus\_intermediate} context
level; Section~\ref{sec:trim} states the executed grid.}
\label{fig:schematic-design}
\end{figure}

\subsection{Unit of observation}

The unit of observation is a \emph{cell}: one complete configuration evaluated on one
benchmark with one seed. A cell is a frozen record of eight factors together with four
fixed parameters, and it is identified by a deterministic string of the form
\code{\{size\}\_\{topology\}[\_n\{k\}]\_\{context\}\_p\{prompt\}\_r\{reasoning\}\_\{benchmark\}\_s\{seed\}},
plus a twelve-character SHA-256 digest of the canonicalised configuration. A cell is
complete when it has written a metadata record and a parseable per-question result
stream for its full question set.

A cell contains roughly 200 graded questions. Within a cell, every question is answered
by the whole system, and every participating agent's per-round output is recorded
separately. The analysis therefore has three grains: cell level, item level (one row per
question) and agent level (one row per agent per round).

A \emph{design cell} is a configuration without its seed. Design cells are the resampling
unit for the second bootstrap stage and the aggregation unit for every marginal mean
reported here, so that a configuration that happened to complete three seeds does not
receive three times the weight of one that completed one.

\subsection{Axes and levels}

Table~\ref{tab:grid} states the executed grid.

\intable{tab_grid}

Each axis was chosen to isolate a lever that is usually confounded in the literature.

\textbf{Capacity} uses six dense Qwen3 checkpoints spanning roughly a factor of fifty in
parameter count \citep{qwen3}. Holding the family, tokenizer and instruction-tuning
recipe fixed means the capacity contrast is not also a training-data or alignment
contrast. Parameter count is a proxy for pretraining compute; the sweep does not observe
training budgets, so capacity claims are stated as parameter-scaling statements.

\textbf{Reasoning budget} varies test-time compute through the model's native thinking
mode rather than through prompt-level instructions to think. The five rungs are a
disabled condition, three finite token budgets spanning a factor of sixteen, and an
unlimited condition. Because the realized number of thinking tokens is recorded per
generation, the reasoning axis can be analysed both as an ordinal knob and as a
continuous dose, and these give different answers when a model does not spend the budget
it is given.

\textbf{Prompt complexity} contrasts a one-line instruction against a six-stage
structured protocol that includes an explicit calibration step. This axis is included
because prompt scaffolding is the cheapest available intervention and because the L3
template instructs the model to report confidence honestly, which makes it a direct test
of whether calibration instructions work.

\textbf{Context sharing} varies what an agent sees of its peers: the final answer alone,
or the answer plus the peer's raw chain of thought. This is the axis that most directly
tests the correlated-error mechanism, since exposing rationales is the mechanism by which
one agent's error can propagate to another.

\textbf{Topology} varies the coordination structure across four levels, including a
single-agent condition that serves as the baseline for every efficiency metric.

\textbf{Benchmark} spans two knowledge-heavy multiple-choice suites of different
difficulty \citep{rein2024gpqa,wang2024mmlupro}, a free-form mathematical suite
\citep{hendrycks2021math,lightman2024verify}, and an adversarial truthfulness suite
\citep{lin2022truthfulqa}. The mathematical suite is included specifically because it has
no option-logprob confidence path, which isolates which conclusions depend on that signal
being available.

\textbf{Seed} replicates each configuration three times. The seed also drives the
deterministic per-question option shuffle, so replication varies both sampling noise and
answer-position effects.

\subsection{Factorial size, collapsing and the trim}
\label{sec:trim}

The raw cross-product of the seven axes is $6 \times 4 \times 2 \times 2 \times 5 \times
4 \times 3 = 5{,}760$ combinations. Two canonicalisation rules reduce this before
execution.

First, the single-agent topology has no peers, so its context-sharing level, agent count
and round count are forced to artifact-only, one and one respectively. Second, the
independent topology has no communication, so its context-sharing level is forced to
artifact-only. Both rules collapse combinations that would otherwise be distinct labels
for identical experiments, and deduplication by cell identifier removes the resulting
repeats. After collapsing, the grid holds 4{,}320 core cells: 720 single-agent, 720
independent, 1{,}440 decentralized and 1{,}440 centralized.

The generator then injects 360 additional single-agent baseline cells, one per
combination of capacity, benchmark, seed and reasoning rung, at prompt level 1. These
exist so that every multi-agent cell has a matched single-agent baseline available for
the efficiency metrics of Section~\ref{sec:efficiency-defs}. The executed manifest is
therefore 4{,}680 cells.

This grid is the result of a trim applied on 2026-06-10, after 1{,}406 of the original
11{,}520 cells had completed. The trim reduced prompt complexity from four levels to two
by keeping the endpoints L0 and L3, and reduced context sharing from three levels to two
by keeping the endpoints artifact-only and plus-CoT. The stated rationale was that in the
partial data both dropped dimensions showed the weakest signal, with an accuracy spread
of 0.035 across the four prompt levels and 0.046 across the three context levels. Of the
1{,}406 already-completed cells, 559 fell inside the new slice and 847 became orphaned;
the orphans were preserved on disk but excluded from the grid. The pre-trim manifest of
11{,}520 cells is retained alongside the current one.

Three consequences of the trim matter for reading the results.

The \code{plus\_intermediate} context level was never executed post-trim, so the
context-sharing variable is effectively binary in the analysis population and the ordinal
context rank takes only the values 0 and 2.

Prompt level 1 is not part of the stated grid but is present in the data through the 360
injected baselines. The core-grid membership predicate requires the prompt level to be 0
or 3, so every injected baseline is excluded from the analysis population even though the
sweep still produces it. This accounts for the single largest block of cells that are
complete but not analysed.

The trim also defines a calendar boundary. A separate flag records whether a cell started
before or after the trim commit epoch. This is distinct from factor-level membership: a
cell can carry post-trim factor levels and a pre-trim timestamp. Both are reported, and
the calendar version is used only as a sensitivity check for serving-stack drift.

\subsection{What the design does and does not identify}

The design identifies the marginal effect of each axis and, within the limits of cell
coverage, two-way interactions among them. It identifies the one-against-three system
contrast. It does not identify an agent-count scaling law, because agent count takes only
two values and the value three is perfectly confounded with having a multi-agent
topology.

Two further confounds are structural rather than accidental. The decentralized and
centralized topologies run two rounds while independent and single-agent run one, so a
topology contrast against those two is also a compute contrast. And single-agent cells
run five self-consistency samples while multi-agent cells do not, so the one-against-three
contrast compares two system designs rather than isolating the number of agents. Both are
stated again in Section~\ref{sec:limitations} and are the reason the agent-count axis is
reported as a design-family comparison throughout.
